\documentclass[10pt,a4paper]{amsart}
\usepackage{amsmath,amssymb,amsthm,graphicx}
% Retained preamble for the Stacks Project chapter snapshot in this directory.
% The chapter's own preamble.tex is a copy of the shared Stacks Project preamble; the block
% below is sliced MECHANICALLY, line for line, from it. Nothing was added, removed or
% redefined; the wrapper class, the xr-hyper/hyperref wiring and the externaldocument list
% are omitted because the assembler supplies the document class.
%
% lines taken from preamble.tex: 3, 5, 9, 11, 12, 156, 157, 158, ... (51 lines total)
%
% package coverage restored from the chapter's own preamble (2026-10-13): xy, xy-2cell

\usepackage[all]{xy}\usepackage{verbatim}
\xyoption{2cell}
\UseAllTwocells

\newenvironment{reference}{\comment}{\endcomment}
\newenvironment{slogan}{\comment}{\endcomment}
\newenvironment{history}{\comment}{\endcomment}
\theoremstyle{plain}
\newtheorem{theorem}[subsection]{Theorem}
\newtheorem{proposition}[subsection]{Proposition}
\newtheorem{lemma}[subsection]{Lemma}
\theoremstyle{definition}
\newtheorem{definition}[subsection]{Definition}
\newtheorem{example}[subsection]{Example}
\newtheorem{exercise}[subsection]{Exercise}
\newtheorem{situation}[subsection]{Situation}
\theoremstyle{remark}
\newtheorem{remark}[subsection]{Remark}
\newtheorem{remarks}[subsection]{Remarks}
\numberwithin{equation}{subsection}
\def\lim{\mathop{\mathrm{lim}}\nolimits}
\def\colim{\mathop{\mathrm{colim}}\nolimits}
\def\Spec{\mathop{\mathrm{Spec}}}
\def\Hom{\mathop{\mathrm{Hom}}\nolimits}
\def\Ext{\mathop{\mathrm{Ext}}\nolimits}
\def\SheafHom{\mathop{\mathcal{H}\!\mathit{om}}\nolimits}
\def\SheafExt{\mathop{\mathcal{E}\!\mathit{xt}}\nolimits}
\def\Sch{\mathit{Sch}}
\def\Mor{\mathop{\mathrm{Mor}}\nolimits}
\def\Ob{\mathop{\mathrm{Ob}}\nolimits}
\def\Sh{\mathop{\mathit{Sh}}\nolimits}
\def\NL{\mathop{N\!L}\nolimits}
\def\CH{\mathop{\mathrm{CH}}\nolimits}
\def\proetale{{pro\text{-}\acute{e}tale}}
\def\etale{{\acute{e}tale}}
\def\QCoh{\mathit{QCoh}}
\def\Ker{\mathop{\mathrm{Ker}}}
\def\Im{\mathop{\mathrm{Im}}}
\def\Coker{\mathop{\mathrm{Coker}}}
\def\Coim{\mathop{\mathrm{Coim}}}
\DeclareMathSymbol{\boxtimes}{\mathbin}{AMSa}{"02}
\def\QCohstack{\mathcal{QC}\!\mathit{oh}}
\def\Cohstack{\mathcal{C}\!\mathit{oh}}
\def\Spacesstack{\mathcal{S}\!\mathit{paces}}
\def\Quotfunctor{\mathrm{Quot}}
\def\Hilbfunctor{\mathrm{Hilb}}
\def\Curvesstack{\mathcal{C}\!\mathit{urves}}
\def\Polarizedstack{\mathcal{P}\!\mathit{olarized}}
\def\Complexesstack{\mathcal{C}\!\mathit{omplexes}}
\def\Pic{\mathop{\mathrm{Pic}}\nolimits}
\def\Picardstack{\mathcal{P}\!\mathit{ic}}
\def\Picardfunctor{\mathrm{Pic}}
\def\Deformationcategory{\mathcal{D}\!\mathit{ef}}
\usepackage{cleveref}
\crefname{originallemma}{original lemma}{original lemmas}
\Crefname{originallemma}{Original lemma}{Original lemmas}
\makeatletter
\newlabel{lemma-affine}{{lemma-affine}{source}{}{sourcepointer-1}{}}
\newlabel{lemma-affine@cref}{{[originallemma][0][]lemma-affine}{[1][0][]source}{}{}{}}
\makeatother
\crefname{originallemma}{original lemma}{original lemmas}
\Crefname{originallemma}{Original lemma}{Original lemmas}
\makeatletter
\newlabel{lemma-easy-kunneth}{{lemma-easy-kunneth}{source}{}{sourcepointer-2}{}}
\newlabel{lemma-easy-kunneth@cref}{{[originallemma][0][]lemma-easy-kunneth}{[1][0][]source}{}{}{}}
\makeatother
\crefname{originallemma}{original lemma}{original lemmas}
\Crefname{originallemma}{Original lemma}{Original lemmas}
\makeatletter
\newlabel{lemma-flat}{{lemma-flat}{source}{}{sourcepointer-3}{}}
\newlabel{lemma-flat@cref}{{[originallemma][0][]lemma-flat}{[1][0][]source}{}{}{}}
\makeatother
\crefname{originallemma}{original lemma}{original lemmas}
\Crefname{originallemma}{Original lemma}{Original lemmas}
\makeatletter
\newlabel{lemma-fully-faithful}{{lemma-fully-faithful}{source}{}{sourcepointer-4}{}}
\newlabel{lemma-fully-faithful@cref}{{[originallemma][0][]lemma-fully-faithful}{[1][0][]source}{}{}{}}
\makeatother
\crefname{originallemma}{original lemma}{original lemmas}
\Crefname{originallemma}{Original lemma}{Original lemmas}
\makeatletter
\newlabel{lemma-universally-injective}{{lemma-universally-injective}{source}{}{sourcepointer-5}{}}
\newlabel{lemma-universally-injective@cref}{{[originallemma][0][]lemma-universally-injective}{[1][0][]source}{}{}{}}
\makeatother
\crefname{originallemma}{original lemma}{original lemmas}
\Crefname{originallemma}{Original lemma}{Original lemmas}
\makeatletter
\newlabel{proposition-stack-fpqc}{{proposition-stack-fpqc}{source}{}{sourcepointer-6}{}}
\newlabel{proposition-stack-fpqc@cref}{{[originallemma][0][]proposition-stack-fpqc}{[1][0][]source}{}{}{}}
\makeatother
\crefname{originallemma}{original lemma}{original lemmas}
\Crefname{originallemma}{Original lemma}{Original lemmas}
\makeatletter
\newlabel{stacks-cohomology-lemma-relative-leray}{{0783}{source}{}{sourcepointer-7}{}}
\newlabel{stacks-cohomology-lemma-relative-leray@cref}{{[originallemma][0][]0783}{[1][0][]source}{}{}{}}
\makeatother
\crefname{originallemma}{original lemma}{original lemmas}
\Crefname{originallemma}{Original lemma}{Original lemmas}
\makeatletter
\newlabel{stacks-properties-lemma-quasi-compact-stack}{{04YC}{source}{}{sourcepointer-8}{}}
\newlabel{stacks-properties-lemma-quasi-compact-stack@cref}{{[originallemma][0][]04YC}{[1][0][]source}{}{}{}}
\makeatother
\title{Functors on QCoh come from morphisms (affine diagonal case)}
\author{The Stacks Project Authors}
\date{stacks-project:stacks-more-morphisms.tex:local-snapshot-20261012; GNU FDL 1.2 or later}
\begin{document}
\maketitle
\noindent\textit{Editorial note.} Statement span is lines 4847-4854 of sources/stacks-stacks-more-morphisms/pinned-original-source.tex; proof span is lines 4856-4930, the complete author proof (73 body lines).\par\medskip
\noindent\textit{Editorial note.} Navigation stubs: 8 cited-but-not-included labels are supplied as non-load-bearing pointers with their pinned tags where the snapshot tag index has them.\par\medskip
\noindent\textit{Editorial note.} extra\_packages is empty on purpose: the retained preamble already loads xy with options and a second declaration clashes in the assembler wrapper.\par\medskip
\noindent\textit{Editorial note.} Text is retained verbatim, including the author's own typos and placeholders.\par\medskip
\noindent\textit{Editorial note.} Queue index 16 of queue\_stacks-stacks-more-morphisms.json (pb=2782, stmt 4847-4854, proof 4856-4930).\par\medskip
\noindent\textit{Editorial note.} stacks-stacks-more-morphisms wave smm01 (2026-10-14): fresh adjudication accept per lane-s3/smm\_judgments.md; pool per Chief 10:5x open-pool ruling (net 59 with-proof rows, bloodline lane-s3 by anchor audit, v2 true-key namespace stacks:stacks-more-morphisms.tex, row lemma-deform-property-fp-over-ft avoided on cross-chapter key collision); 60 with-proof rows judged 45 accept / 14 reject / 1 avoid; patch-sd-envs contract spans scanned. Named omission: top-square commutativity and the functor correspondence 'details omitted' - named-mechanism bookkeeping; uniqueness rides on the fully-faithful lemma.\par\medskip
\noindent\textit{Allowed standard prerequisites (editorial).}\begin{itemize}
\item Explicit prerequisite boundary. The cited inputs are prerequisites; the argument in the retained span is not replaced by any cited result.
\end{itemize}\medskip
\ifdefined\hypertarget\hypertarget{sourcepointer-1}{}\fi
\noindent\textit{Source reference.} Original lemma lemma-affine: sources/stacks-stacks-more-morphisms/pinned-original-source.tex, source lines 4575--4583. Cited by the retained statement or proof but not part of the collected spans; supplied only as a navigation pointer. The argument does not depend on its content.\par\medskip
\ifdefined\hypertarget\hypertarget{sourcepointer-2}{}\fi
\noindent\textit{Source reference.} Original lemma lemma-easy-kunneth: sources/stacks-stacks-more-morphisms/pinned-original-source.tex, source lines 4738--4747. Cited by the retained statement or proof but not part of the collected spans; supplied only as a navigation pointer. The argument does not depend on its content.\par\medskip
\ifdefined\hypertarget\hypertarget{sourcepointer-3}{}\fi
\noindent\textit{Source reference.} Original lemma lemma-flat: sources/stacks-stacks-more-morphisms/pinned-original-source.tex, source lines 4716--4721. Cited by the retained statement or proof but not part of the collected spans; supplied only as a navigation pointer. The argument does not depend on its content.\par\medskip
\ifdefined\hypertarget\hypertarget{sourcepointer-4}{}\fi
\noindent\textit{Source reference.} Original lemma lemma-fully-faithful: sources/stacks-stacks-more-morphisms/pinned-original-source.tex, source lines 4792--4799. Cited by the retained statement or proof but not part of the collected spans; supplied only as a navigation pointer. The argument does not depend on its content.\par\medskip
\ifdefined\hypertarget\hypertarget{sourcepointer-5}{}\fi
\noindent\textit{Source reference.} Original lemma lemma-universally-injective: sources/stacks-stacks-more-morphisms/pinned-original-source.tex, source lines 4646--4650. Cited by the retained statement or proof but not part of the collected spans; supplied only as a navigation pointer. The argument does not depend on its content.\par\medskip
\ifdefined\hypertarget\hypertarget{sourcepointer-6}{}\fi
\noindent\textit{Source reference.} Original lemma proposition-stack-fpqc: sources/stacks-stacks-more-morphisms/pinned-original-source.tex, source lines 4321--4325. Cited by the retained statement or proof but not part of the collected spans; supplied only as a navigation pointer. The argument does not depend on its content.\par\medskip
\ifdefined\hypertarget\hypertarget{sourcepointer-7}{}\fi
\noindent\textit{Source reference.} Original lemma 0783: sources/stacks-stacks-more-morphisms/pinned-original-source.tex, source lines 4930--4930. Cited by the retained statement or proof but not part of the collected spans; supplied only as a navigation pointer. The argument does not depend on its content.\par\medskip
\ifdefined\hypertarget\hypertarget{sourcepointer-8}{}\fi
\noindent\textit{Source reference.} Original lemma 04YC: sources/stacks-stacks-more-morphisms/pinned-original-source.tex, source lines 4930--4930. Cited by the retained statement or proof but not part of the collected spans; supplied only as a navigation pointer. The argument does not depend on its content.\par\medskip
\begin{lemma}
\label{lemma-main}
Let $\mathcal{X}$ be a Noetherian algebraic stack with affine diagonal.
Let $B$ be a ring.
Let $F : \QCoh(\mathcal{O}_\mathcal{X}) \to \text{Mod}_B$
be a right exact tensor functor which commutes with direct sums.
Then $F$ comes from a unique morphism $\Spec(B) \to \mathcal{X}$.
\end{lemma}

\begin{proof}
Choose a surjective smooth morphism $g : U \to \mathcal{X}$ with $U = \Spec(A)$
affine, see Properties of Stacks, Lemma
\ref{stacks-properties-lemma-quasi-compact-stack}. Apply
Lemma \ref{lemma-affine} to get the finite type
commutative $B$-algebra $C = F(g_{\QCoh, *}\mathcal{O}_U)$
and the ring map $A \to C$.
By Lemma \ref{lemma-universally-injective}
the ring map $B \to C$ is universally injective.
Consider the algebra
$$
C \otimes_B C =
F(g_{\QCoh, *}\mathcal{O}_U
\otimes_{\mathcal{O}_\mathcal{X}}
g_{\QCoh, *}\mathcal{O}_U)
$$
Since $g$ is flat, quasi-compact, and quasi-separated
by Lemma \ref{lemma-easy-kunneth} we have the first equality in
$$
g_{\QCoh, *}\mathcal{O}_U
\otimes_{\mathcal{O}_\mathcal{X}}
g_{\QCoh, *}\mathcal{O}_U
=
f_{\QCoh, *}\mathcal{O}_{U \times_\mathcal{X} U} =
g_{\QCoh, *}(\text{pr}_{2, *}\mathcal{O}_{U \times_\mathcal{X} U})
$$
where
$f : U \times_\mathcal{X} U \to \mathcal{X}$ is the obvious morphism
and $\text{pr}_2 : U \times_\mathcal{X} U \to U$ is the second
projection. The second equality follows from
Cohomology of Stacks, Lemma \ref{stacks-cohomology-lemma-relative-leray}
and $f = g \circ \text{pr}_2$. Since the diagonal of $\mathcal{X}$
is affine, we see that $U \times_\mathcal{X} U = \Spec(R)$ is affine.
Let us use $\text{pr}_2 : A \to R$ to view $R$ as an $A$-algebra.
All in all we obtain
$$
C \otimes_B C =
F(g_{\QCoh, *}\mathcal{O}_U
\otimes_{\mathcal{O}_\mathcal{X}}
g_{\QCoh, *}\mathcal{O}_U) =
F(g_{\QCoh, *}(\text{pr}_{2, *}\mathcal{O}_{U \times_\mathcal{X} U})) =
R \otimes_A C
$$
where the final equality follows from the final statement of
Lemma \ref{lemma-affine}. Since $A \to R$ is flat (because
$\text{pr}_2$ is flat as a base change of $U \to \mathcal{X}$),
we conclude that $C \otimes_B C$ is flat over $C$.
By Lemma \ref{lemma-flat} we conclude that $B \to C$ is faithfully flat.

\medskip\noindent
We claim there is a solid commutative diagram
$$
\xymatrix{
\Spec(C \otimes_B C) \ar@<1ex>[d] \ar@<-1ex>[d] \ar[r] &
U \times_\mathcal{X} U \ar@<1ex>[d] \ar@<-1ex>[d] \\
\Spec(C) \ar[d] \ar[r] &
U \ar[d] \\
\Spec(B) \ar@{..>}[r] &
\mathcal{X}
}
$$
The arrow $\Spec(C) \to U = \Spec(A)$ comes from the ring map
$A \to C$ in the statement of Lemma \ref{lemma-affine}.
The arrow $\Spec(C \otimes_B C) \to U \times_\mathcal{X} U$
similarly comes from the ring map $R \to C \otimes_B C$.
To verify the top square commutes use Lemma \ref{lemma-fully-faithful};
details omitted.
We conclude we get the dotted arrow $\Spec(B) \to \mathcal{X}$
by Proposition \ref{proposition-stack-fpqc}.

\medskip\noindent
The statement that $F$ is the functor corresponding to pullback
by the dotted arrow is also clear from this and the corresponding
statement in Lemma \ref{lemma-affine}. Details omitted.
\end{proof}

\end{document}
